El presente trabajo constituye una contribución original al campo de la ingeniería de software al integrar tres áreas de investigación de frontera: los modelos de lenguaje de gran escala, las arquitecturas multi-agente y el MCP, aplicados a la automatización de un proceso tan fundamental como la especificación de requisitos de software. Académicamente, el sistema propuesto aporta un marco de referencia implementable para futuras investigaciones que busquen extender la automatización hacia otras fases del ciclo de vida del desarrollo de software. Asimismo, la metodología de evaluación basada en juicio de expertos con criterios derivados del estándar IEEE 830-1998 y la metodología de elicitación establece un precedente metodológico reproducible para trabajos similares en el ámbito hispanohablante, donde la documentación técnica en español representa un desafío adicional para los sistemas de procesamiento de lenguaje natural. 


Desde la perspectiva investigativa, este trabajo responde directamente a las brechas identificadas en la literatura científica reciente. Wolf et al. \cite{Wolf2026} señalan que el campo de la ingeniería de requisitos asistida por IA carece de integración con procesos reales de ingeniería de software y de estandarización metodológica. El sistema propuesto aborda precisamente esta brecha al conectar el proceso real de elicitación —la entrevista de audio— con la generación formal del documento SRS bajo un estándar reconocido internacionalmente. En este contexto, investigaciones previas desarrolladas en la Universidad Técnica del Norte han demostrado la viabilidad de aplicar LLMs en procesos de análisis automatizado de documentos técnicos \cite{Landeta-Lpez2024}, \cite{Guevara-Vega2019} validando la pertinencia de esta línea de investigación en el contexto institucional y estableciendo un antecedente directo sobre el que la presente propuesta busca avanzar. Adicionalmente, la incorporación del MCP como infraestructura de comunicación entre agentes representa una línea de investigación emergente con escasa exploración en el contexto de la ingeniería de requisitos, lo que posiciona este trabajo en la frontera del conocimiento actual en sistemas multiagente basados en LLMs\cite{Guo2024}, \cite{Ray2025}.
Técnicamente, el sistema aportará una solución concreta y funcional para uno de los problemas más costosos del desarrollo de software: la elaboración manual del documento SRS. La arquitectura multi-agente propuesta permite distribuir la complejidad del proceso en agentes especializados, superando las limitaciones de los modelos de lenguaje individuales en términos de gestión de contextos extensos y propensión a generar alucinaciones. \cite{Rahman2024}. La integración del agente de búsqueda web mediante RAG (Retrieval-Augmented Generation) enriquecerá el contexto del dominio específico del software a desarrollar, mejorando la pertinencia y completitud de los requisitos generados \cite{Mukaida2025}. El sistema también aportará un repositorio persistente de documentos SRS que permitirá la consulta, comparación y reutilización de requisitos entre proyectos, constituyendo una base de conocimiento organizacional de alto valor para equipos de desarrollo. 
Desde el punto de vista económico, la automatización de la generación del documento SRS representa un ahorro significativo en horas de trabajo de los ingenieros más experimentados del equipo, quienes actualmente destinan una parte considerable de su tiempo a tareas de transcripción, clasificación y redacción formal de requisitos \cite{Rahman2024}.  

Este tiempo liberado puede ser redirigido hacia actividades de mayor valor como el diseño de arquitectura, la revisión de código o la gestión de riesgos del proyecto. Para organizaciones de desarrollo de software de tamaño pequeño y mediano, donde los recursos humanos son limitados, esta reducción en el esfuerzo de documentación puede representar una ventaja competitiva significativa y una reducción directa en los costos asociados a retrabajos provocados por especificaciones deficientes \cite{Wolf2026}, \cite{Umar2024}. 

Socialmente, el sistema beneficiará a múltiples actores del ecosistema de desarrollo de software. Los ingenieros de software verán reducida su carga operativa en tareas repetitivas y podrán enfocarse en actividades que requieren mayor creatividad y juicio profesional. Los usuarios finales y clientes se beneficiarán de documentos de requisitos más completos y consistentes, lo que reduce la probabilidad de que el sistema construido no satisfaga sus necesidades reales \cite{ieee1984ieee}. En el ámbito educativo, el sistema tiene un potencial pedagógico significativo: estudiantes de ingeniería de software podrán utilizarlo como herramienta de apoyo para comprender la estructura y contenido de un documento SRS conforme al estándar IEEE 830-1998, aprendiendo de manera práctica los elementos que debe contener una especificación de calidad \cite{doe2011recommended}, \cite{ieee1984ieee}. La metodología de elicitación de Durán y Bernárdez \cite{toro2000metodologia}, ampliamente utilizada en la enseñanza de ingeniería de requisitos en el mundo hispanohablante, servirá como referente pedagógico integrado en el sistema. 

En síntesis, el sistema propuesto resolverá un problema concreto y medible: la brecha existente entre el lenguaje natural de las entrevistas de elicitación y la especificación técnica formal requerida por el estándar IEEE 830-1998. Sus beneficiarios directos son los ingenieros de software, equipos de desarrollo, organizaciones de tecnología y estudiantes de ingeniería. Sus aportes técnicos incluyen una arquitectura multi-agente replicable, un flujo automatizado de procesamiento desde el audio de la entrevista hasta la generación del documento SRS, y un conjunto de prompts estructurados para la extracción y clasificación de requisitos. Sus aportes teóricos incluyen evidencia empírica sobre la viabilidad del uso de LLMs multi-agente con MCP en la automatización de procesos de ingeniería de requisitos, contribuyendo al avance del conocimiento en un área de alta relevancia para la industria del software a nivel global.